% (a) the scan graph with weights; (b) physics parallel: double well, ground state and slowest mode.
\begin{tikzpicture}[fig]
\begin{scope}
  \node[lab] at (2.0,3.9) {(a) weights $w=e^{-(J/\tilde J)^2}$ on the links};
  % links: thick inside each side, thin across the wall between columns 2 and 3
  \foreach \j in {0,...,3} {
    \foreach \i in {0,1} \draw[line width=2.2pt, cP0!45] (\i,\j) -- (\i+1,\j);
    \foreach \i in {3} \draw[line width=2.2pt, cP1!45] (\i,\j) -- (\i+1,\j);
    \draw[line width=0.3pt, cInk2] (2,\j) -- (3,\j);
  }
  \foreach \i in {0,1,2} \foreach \j in {0,1,2} \draw[line width=2.2pt, cP0!45] (\i,\j) -- (\i,\j+1);
  \foreach \i in {3,4} \foreach \j in {0,1,2} \draw[line width=2.2pt, cP1!45] (\i,\j) -- (\i,\j+1);
  \foreach \i in {0,1,2} \foreach \j in {0,...,3} \node[dot=cP0, minimum size=6pt] at (\i,\j) {};
  \foreach \i in {3,4} \foreach \j in {0,...,3} \node[dot=cP1, minimum size=6pt] at (\i,\j) {};
  \draw[wall, dashed] (2.5,-0.4) -- (2.5,3.4);
  \node[small, text width=4.6cm] at (2,-0.95) {thick: small jump, easy to cross;
    thin: large jump. The walker stays on one side for a long time.};
\end{scope}
\begin{scope}[shift={(6.6,0.2)}]
  \node[lab] at (2.2,3.1) {(b) physics parallel: two basins};
  \draw[->] (0,0) -- (4.6,0) node[right, small] {$x$};
  \draw[thick, cInk2, domain=0.2:4.2, samples=80, smooth] plot (\x, {0.35*((\x-2.2)^2-1.3)^2/1.7+0.25});
  \node[small, cInk2] at (0.05,2.2) {$U(x)$};
  \draw[cP2, thick, domain=0.2:4.2, samples=60, smooth] plot (\x, {1.25+0.0*\x});
  \node[small, cP2, anchor=west] at (4.25,1.25) {$\psi_0$: equilibrium};
  \draw[cAcc, thick, domain=0.2:4.2, samples=80, smooth] plot (\x, {1.25+0.5*tanh(3*(\x-2.2))});
  \node[small, cAcc, anchor=west] at (4.25,1.75) {$\psi_1$: slowest mode};
  \node[small, text width=4.8cm] at (2.2,-0.75) {$\psi_1\approx\pm1$ in each basin and changes sign at the barrier;
    its rate $\lambda_1$ is the escape rate (Kramers).};
\end{scope}
\end{tikzpicture}
